% Options for packages loaded elsewhere
\PassOptionsToPackage{unicode}{hyperref}
\PassOptionsToPackage{hyphens}{url}
\documentclass[
]{article}
\usepackage{xcolor}
\usepackage{amsmath,amssymb}
\setcounter{secnumdepth}{-\maxdimen} % remove section numbering
\usepackage{iftex}
\ifPDFTeX
  \usepackage[T1]{fontenc}
  \usepackage[utf8]{inputenc}
  \usepackage{textcomp} % provide euro and other symbols
\else % if luatex or xetex
  \usepackage{unicode-math} % this also loads fontspec
  \defaultfontfeatures{Scale=MatchLowercase}
  \defaultfontfeatures[\rmfamily]{Ligatures=TeX,Scale=1}
\fi
\usepackage{lmodern}
\ifPDFTeX\else
  % xetex/luatex font selection
\fi
% Use upquote if available, for straight quotes in verbatim environments
\IfFileExists{upquote.sty}{\usepackage{upquote}}{}
\IfFileExists{microtype.sty}{% use microtype if available
  \usepackage[]{microtype}
  \UseMicrotypeSet[protrusion]{basicmath} % disable protrusion for tt fonts
}{}
\makeatletter
\@ifundefined{KOMAClassName}{% if non-KOMA class
  \IfFileExists{parskip.sty}{%
    \usepackage{parskip}
  }{% else
    \setlength{\parindent}{0pt}
    \setlength{\parskip}{6pt plus 2pt minus 1pt}}
}{% if KOMA class
  \KOMAoptions{parskip=half}}
\makeatother
\usepackage{longtable,booktabs,array}
\newcounter{none} % for unnumbered tables
\usepackage{calc} % for calculating minipage widths
% Correct order of tables after \paragraph or \subparagraph
\usepackage{etoolbox}
\makeatletter
\patchcmd\longtable{\par}{\if@noskipsec\mbox{}\fi\par}{}{}
\makeatother
% Allow footnotes in longtable head/foot
\IfFileExists{footnotehyper.sty}{\usepackage{footnotehyper}}{\usepackage{footnote}}
\makesavenoteenv{longtable}
\usepackage{graphicx}
\makeatletter
\newsavebox\pandoc@box
\newcommand*\pandocbounded[1]{% scales image to fit in text height/width
  \sbox\pandoc@box{#1}%
  \Gscale@div\@tempa{\textheight}{\dimexpr\ht\pandoc@box+\dp\pandoc@box\relax}%
  \Gscale@div\@tempb{\linewidth}{\wd\pandoc@box}%
  \ifdim\@tempb\p@<\@tempa\p@\let\@tempa\@tempb\fi% select the smaller of both
  \ifdim\@tempa\p@<\p@\scalebox{\@tempa}{\usebox\pandoc@box}%
  \else\usebox{\pandoc@box}%
  \fi%
}
% Set default figure placement to htbp
\def\fps@figure{htbp}
\makeatother
\setlength{\emergencystretch}{3em} % prevent overfull lines
\providecommand{\tightlist}{%
  \setlength{\itemsep}{0pt}\setlength{\parskip}{0pt}}
\usepackage{bookmark}
\IfFileExists{xurl.sty}{\usepackage{xurl}}{} % add URL line breaks if available
\urlstyle{same}
% fallback for those not using the hyperref driver hyperxmp:
\makeatletter
\@ifundefined{xmpquote}{\newcommand{\xmpquote}[1]{#1}}{}
\makeatother
\hypersetup{
  pdftitle={Auditing constrained geospatial allocation under noisy annual land-cover labels},
  hidelinks,
  pdfcreator={LaTeX via pandoc}}

\title{Auditing constrained geospatial allocation under noisy annual
land-cover labels}
\author{}
\date{}

\begin{document}
\maketitle

\subsection{Abstract}\label{abstract}

Urban land-change models are often evaluated against annual labels whose
reliability is not spatially uniform. We present an auditable protocol
for testing constrained spatial allocation under that uncertainty, using
Geospatial Kernel, the algorithmic core of a Geospatial World Model, as
the focal execution layer. The Abu Dhabi public-data benchmark uses
annual 100-m land-cover products from 2017--2024, three computational
seeds, destination-correct multi-class change Figure of Merit (FoM),
spatial-block bootstrap contrasts and explicit state--action--constraint
traces. In the primary unmatched-pipeline test, strict FoM was 0.1961
for Geospatial Kernel, 0.1256 for a GeoSOS-derived FLUS-style ANN--CA
console and 0.1619 for GeoFM-LDN in 2023; GeoFM-LDN led after recursive
two-step rollout in 2024. A separate leakage-controlled expanding-window
analysis, excluding the later road snapshot, gave Kernel FoM values of
0.2595, 0.0724, 0.2116 and 0.1816 for the four one-step targets from
2021 to 2024, exceeding a random minimum-change allocator in every
window. In the only available WorldCover 2020--2021 window, Dynamic
World built pixels decreased, so the oracle-demand built-gain test had
zero power for every allocator. Dynamic World-to-WorldCover built-gain
F1 was 0.0111--0.0192. The count-controlled allocation variant was near
opportunity-level overall: Kernel was below random at 0.05,
indistinguishable from random at 0.10--0.20 given three-seed variation,
and above random only at 0.30. The planning track is therefore reported
as scenario stress testing, not official forecasting. The protocol's
principal contribution is an inspectable proposal--projection--writeback
boundary; the experiments support conditional allocation auditability,
not verified prediction of Abu Dhabi's actual future.

\subsection{Keywords}\label{keywords}

Geospatial World Model; land-cover change; constrained allocation;
cellular automata; Abu Dhabi; scenario planning

\subsection{Introduction}\label{introduction}

Land-cover change is a spatial process. Development demand is expressed
as a quantity, but its consequences depend on where each new cell is
placed relative to roads, existing built areas, ecological features and
other land-cover classes. This coupling matters for urban planning
because two maps with the same total built area can imply very different
infrastructure requirements and environmental pressures.
Cellular-automata and suitability-based models have long provided a way
to distribute land demand across a map, including CLUE-S, SLEUTH,
Dinamica EGO, FLUS, local-attribute urban CA models and more recent
patch-generating or deep-learning CA models (Verburg et al., 2002; Silva
and Clarke, 2002; Soares-Filho et al., 2002; Lau and Kam, 2005; Liu et
al., 2017; Liang et al., 2021). Global Earth-observation products now
make it possible to construct annual, spatially aligned inputs for
cities that lack a harmonized local data archive (Brown et al., 2022;
Gorelick et al., 2017).

Validation studies have also shown that agreement in unchanged cells can
mask poor allocation of actual change. Quantity and allocation
disagreement, destination-specific transition scores and explicit
calibration/validation protocols are therefore preferable to a single
overall-accuracy number (Pontius et al., 2008; Pontius and Millones,
2011; van Vliet et al., 2011, 2016). This distinction is central here
because the benchmark uses noisy annual labels and tests both one-step
allocation and recursive rollout.

The main technical difficulty is not producing another categorical
raster. It is keeping the semantics of a planning action separate from
the learned transition tendency. A model may predict that a cell is
likely to become built, yet the resulting map can still violate a water
exclusion, miss the requested class totals or become inconsistent when
the next year is simulated from the model's own output. These failure
modes are particularly important when the model is used to compare
scenarios rather than to reproduce one observed transition.
Latent-dynamics approaches such as GeoFM-LDN can represent continuous
geospatial state and condition transitions on demand, whereas
traditional ANN-plus-cellular-automata systems combine suitability with
explicit neighbourhood evolution. Neither design, by itself, guarantees
that action semantics, hard constraints, state persistence and evidence
provenance are visible at the same execution boundary.

We address this systems problem by studying Geospatial Kernel, the core
transition layer of a Geospatial World Model (GWM). The Kernel
represents one simulation step as a typed state, a typed action, a
transition proposal, a constraint projection and a state writeback. The
learned proposal remains domain-specific; the runtime contract is
domain-neutral. For Abu Dhabi land-cover planning, the proposal is a
six-class probability cube, and the projection allocates only mutable
cells until the action's feasible class totals are reached. This
separation makes the planning operation inspectable: a reviewer can
distinguish what the transition model preferred from what the
constraints admitted.

We evaluate this design in a controlled public-data benchmark rather
than treating the benchmark as an official data release. All candidates
consume the same 100-m grid, common valid mask, annual class-count
actions, hard exclusions and evaluator. The historical track uses
observed target counts to isolate spatial allocation skill from demand
error. The planning track starts from the observed 2024 state and
applies three planner-supplied actions: moderate growth, green-priority
growth and high outward growth. Abu Dhabi Plan 2030 provides an
institutional planning context for discussing urban structure, but the
numerical actions in this benchmark are not calibrated to that plan or
to an official population forecast (Abu Dhabi Urban Planning Council,
2007). We ask four questions: whether the Kernel execution contract is
auditable, whether its allocation skill is competitive, whether it
produces feasible and spatially distinct allocations and which
conclusions survive label and parameter sensitivity.

The paper makes three bounded contributions. First, it gives a concrete
state--action--constraint formulation for using a geospatial kernel as
an algorithmic core rather than as an opaque end-to-end predictor.
Second, it defines an evaluation protocol that combines
confidence-filter retention, expanding-window temporal tests, zero-model
controls and external-product disagreement rather than relying on one
favourable split. Third, it turns future maps into explicit stress-test
artefacts, including cumulative rasters and vectorized transition
footprints, while preserving the distinction between public-data land
cover and authoritative legal land use.

\subsection{Methods}\label{methods}

\subsubsection{Portable audit protocol}\label{portable-audit-protocol}

The benchmark separates a reusable audit protocol from the Abu Dhabi
adapter. For any city or geospatial domain, the protocol requires: (i) a
declared class crosswalk and spatially stratified label-confidence
tiers; (ii) persistence and minimum-change zero models; (iii)
destination-correct multi-class change FoM with quantity and allocation
diagnostics; (iv) spatial-block bootstrap intervals and paired model
contrasts on a frozen evaluation mask; (v) an explicit check for
structural-zero metrics before Pareto selection; (vi) a versioned
objective table with direction, units and interpretation; and (vii) a
state--action--constraint trace recording proposal, projection,
hard-mask status and state writeback. The protocol also requires every
model/year/seed record to report the number of predicted changed cells.
A record with \texttt{predicted\_change\_pixels\ ==\ 0} is marked as a
degenerate output and reviewed alongside structural-zero metrics; it
cannot by itself be interpreted as model superiority or real land-cover
stability. The protocol reports domain validity separately from
execution completeness and requires independent change validation before
a categorical map is presented as an observed land-use product. A
structural-zero metric is flagged operationally when its denominator is
zero for every candidate under the declared action and constraint rules;
for exact-count allocation, a conversion rate is also a structural zero
when the protected/source rule makes the relevant source class
monotonically non-decreasing, so no eligible source cells can contribute
to that transition. Such a metric is reported descriptively and excluded
from Pareto selection.

\subsubsection{Abu Dhabi implementation: boundary, grid and temporal
protocol}\label{abu-dhabi-implementation-boundary-grid-and-temporal-protocol}

The study boundary is the polygon returned for OpenStreetMap relation
R4479763, representing Abu Dhabi city (OpenStreetMap contributors,
2026). All rasters were aligned to a canonical 100-m grid in EPSG:32640
with 475 columns and 360 rows. A cell was included when its centre fell
within the city polygon and when the common public-data mask indicated
valid coverage. The benchmark profile records the source geometry,
raster transform, dimensions and SHA-256 hashes. The effective
evaluation mask contains 79,726 cells.

The observed sequence contains annual states for 2017--2024. Models were
fit using the four transitions 2017→2018, 2018→2019, 2019→2020 and
2020→2021. A 2021→2022 transition was used for validation during model
development. The historical test starts from 2022 and evaluates 2023 and
2024 in an open-loop rollout without observed-state writeback. The
planning test starts from the observed 2024 state and recursively
generates 2025--2031 under each scenario action.

An additional Geospatial Kernel analysis used four expanding one-step
windows with origins in 2020--2023. For origin year \(t\), training
included only transitions whose target year was at or before \(t\); the
held-out \(t+1\) label was prohibited from feature construction and
fitting. The held-out target supplied its observed class counts to the
oracle-demand allocator and was then used for evaluation. This design
isolates spatial allocation conditional on known demand and is not an
end-to-end demand forecast. Early folds locked only origin-year water
and wetland cells and omitted the two road-distance features, because
applying the 2021 WorldCover proxy or the later OpenStreetMap snapshot
to earlier folds would introduce future information. The resulting
rolling proposal has 23 features. A machine-tested temporal firewall
records the last training target, prediction target and permitted
target-year uses for every run. The headline 2023--2024 analysis is less
strict: it uses the OSM road snapshot accessed on 31 July 2026, which
may include roads built after the corresponding prediction years
(including after 2022). This snapshot is supplied symmetrically to all
three pipelines, so it does not change their ordering, but it can affect
absolute scores; the rolling-origin results are the strict
no-road-snapshot boundary.

\subsubsection{Public data and class
semantics}\label{public-data-and-class-semantics}

Dynamic World V1 provides annual scene-argmax labels and mean top-class
probabilities (Brown et al., 2022). The nine Dynamic World labels were
crosswalked to six canonical classes: water (source 0), woody vegetation
(1), low vegetation (2, 4 and 5), wetland (3), built (6) and bare (7).
Source class 8 was excluded. The resulting label is explicitly treated
as land cover. No residential, commercial or industrial category is
inferred from the raster labels.

Annual 64-dimensional AlphaEarth embeddings were spatially averaged to
the canonical grid and locally L2-normalized (Brown et al., 2025).
Monthly \texttt{NOAA/VIIRS/DNB/MONTHLY\_V1/VCMSLCFG} night-time-light
radiance was averaged by year (Elvidge et al., 2021). The Copernicus DEM
GLO-30 \texttt{COPERNICUS/DEM/GLO30\_2024\_1} collection supplied
elevation and a finite-difference slope derivative (European Space
Agency, 2024). OpenStreetMap road geometries were rasterized into
distance-to-road and distance-to-major-road surfaces. ESA WorldCover
2021 supplied public water and wetland/mangrove constraint proxies
(Zanaga et al., 2022). All dynamic and static layers were checked for
exact CRS, transform, width and height agreement.

For an external public-product diagnostic, the ESA WorldCover 2020 v1.0
and 2021 v2.0 built-up class (value 50) was converted to built fraction
on each 100-m benchmark cell (Zanaga et al., 2021, 2022). Built presence
was tested at fractions of 0.05, 0.10, 0.20 and 0.30. Built gain
required a cell to cross from below to at or above the threshold; built
loss used the reverse definition. Binary precision, recall, F1 and
intersection over union were calculated for the Dynamic World transition
and for the 2020-origin predictions. The original oracle-demand
comparison uses Dynamic World 2021 built counts; because the built count
decreases in this window, it cannot allocate any built gains and is a
structural-zero, zero-power test of gain location. We therefore also
supplied the WorldCover built-gain count as the action, ranked candidate
cells with the 2020-origin Kernel proposal, and compared it with random
allocation from the same candidate population. This count-controlled
variant is the informative allocation diagnostic, but it remains product
agreement rather than ground-truth validation: the two WorldCover years
use different product versions, WorldCover and Dynamic World both derive
from Sentinel observations, and neither represents statutory Abu Dhabi
land use.

\subsubsection{Geospatial Kernel state, action and transition
proposal}\label{geospatial-kernel-state-action-and-transition-proposal}

Let \(S_t(i)\) denote the six-class state at valid cell \(i\), \(A_t\)
the action for the interval \(t\rightarrow t+1\), and \(X_t\) the
aligned driver layers. The Kernel separates a proposal from a
projection:

\[
Q_t(i,c)=p_\theta\big(c\mid x_i(S_t,X_t)\big),
\qquad
S_{t+1}=\Pi_{\mathcal C,A_t}(Q_t,S_t).
\]

The feature vector \(x_i\) contains 25 values: six one-hot current-class
indicators; twelve neighbourhood proportions, comprising 3 × 3 and 7 × 7
windows for each of the six classes; and seven continuous variables,
namely normalized row and column coordinates, clipped elevation, clipped
slope, log-transformed VIIRS radiance, log distance to roads and log
distance to major roads. A scikit-learn histogram gradient-boosting
classifier estimates the six-class probability vector (Pedregosa et al.,
2011). It uses learning rate 0.08, 120 iterations, at most 31 leaf
nodes, minimum leaf size 40 and L2 regularization 1.0.

Training samples combine all changed valid cells and a 20\% random
sample of stable valid cells in each fit transition. The sample weight
is the product of a change emphasis, \(1+5\mathbf{1}(S_t\ne S_{t+1})\),
and a confidence factor, \(0.5+\min(r_t,r_{t+1})\), where \(r\) is the
Dynamic World mean top probability. The target years 2023 and 2024 are
not used during fitting or model selection. Across the three seeds, each
fitted model used 71,892--72,202 pixel rows and all six classes.

\subsubsection{Constraint projection and state
writeback}\label{constraint-projection-and-state-writeback}

For a mutable source cell \(i\) of class \(s\) and target class \(c\),
the base allocation score is

\[
q_{i,s\rightarrow c}=\log p_\theta(c\mid x_i)-\log p_\theta(s\mid x_i)+\lambda\rho^{(7)}_{i,c},
\]

where \(\rho^{(7)}_{i,c}\) is the 7 × 7 fraction of neighbouring cells
currently in class \(c\), and \(\lambda=0.35\). The algorithm computes
class deficits and excesses from the action's feasible target counts,
ranks all admissible source--target candidates by \(q\), and changes
cells in descending order until all deficits and excesses are zero.
Permanent water and the protected wetland/ecological mask are held
fixed. If the counts cannot be satisfied without changing a
hard-exclusion cell, the step fails closed rather than silently
violating the action.

The projected raster is written as the next \texttt{KernelState}. Every
step records a state reference, action evidence, model version,
parameter reference, projection status and diagnostic counts. The
runtime checks domain identity, source-time consistency and action time
advancement. This execution contract is the Geospatial Kernel's
algorithmic contribution; the Abu Dhabi adapter supplies the
land-cover-specific features, learner and constraint masks.

\subsubsection{Baselines and common
evaluator}\label{baselines-and-common-evaluator}

The baseline was run through a GeoSOS-derived FLUS-style ANN--CA console
(modified build; \texttt{paper-benchmark-flus-v1.1}) using seven
continuous drivers: normalized row and column coordinates, elevation,
slope, VIIRS radiance, distance to roads and distance to major roads.
Its ANN suitability surface was passed to a cellular-automata allocation
stage with the common class totals and public hard mask. The archived
run used one 2021 training year, eight ANN hidden neurons, a 3 × 3 CA
neighbourhood, neighbourhood strength 1.0, acceleration factor 0.1 and
1,000 iterations. Its frozen transition matrix protected water and
wetland but did not prohibit built-to-non-built transitions. The
algorithmic base is traceable to the GeoSOS team's public FLUS source
release (Liu et al., 2017; GeoSOS team, n.d.); the accompanying source
repository an archived source repository at an anonymized revision
contains the author modifications used for this paper. Those
modifications add \texttt{train}/\texttt{train-update} entry points and
\texttt{FLUS\_RANDOM\_SEED} deterministic seeding for both ANN sampling
and CA roulette draws, so the build is not the unmodified upstream
executable. Version v1.1 changes only failure handling for invalid input
reads; valid-input algorithmic paths and benchmark rasters are unchanged
from v1. A 25-feature Kernel-matched input was then executed with
absolute paths for seeds 31, 47 and 73 and is reported only as a
platform-sensitive diagnostic. It shares the Kernel feature family, but
not the learning target or projection semantics: the FLUS ANN estimates
same-year label suitability \(p(S_t\mid X_t)\), whereas the Kernel
proposal is trained to estimate a next-state transition
\(p(S_{t+1}\mid S_t,X_t)\) before constrained projection. Matching
features therefore does not match the training task. Two intermediate
configurations were retained as mechanism diagnostics: 13 features
(seven drivers plus current-class indicators) and 19 features (seven
drivers plus neighbourhood fractions). The 25-feature run collapsed to
zero change for macOS seeds 47 and 73 and for all three independent
external Windows x86\_64 verification runs; the remaining macOS seed-31
run is not used as a valid estimator. The external evidence is
identified in Supplementary Table S3; its source rasters are not locally
archived. The 19-feature mode produced changes for all three macOS seeds
but underfilled the observed 4,500 and 8,464 changed cells in 2023 and
2024. Both modes are therefore diagnostic-only.

GeoFM-LDN (Geospatial Foundation-Model Latent Dynamics Network; the
repository's former \texttt{paper58} path identifier) was implemented as
a demand-conditioned residual latent-dynamics network on 64-channel
AlphaEarth patches. Three 128-channel convolutions use dilation rates 1,
2 and 4, followed by a 1 × 1 projection to a 64-dimensional latent
state. A 12-dimensional action vector (origin and target counts for six
classes) is encoded by a two-layer MLP and broadcast across the patch
before concatenation with the embedding. Group normalization and GELU
activations are used in the residual blocks. A scikit-learn
logistic-regression decoder maps the predicted latent embedding to six
semantic classes. Training uses 64 × 64 patches, batch size 2, eight
epochs, AdamW with learning rate \(3\times10^{-4}\) and weight decay
\(10^{-4}\); the loss is the weighted sum of cosine embedding loss, 0.5
semantic cross-entropy, 0.1 demand-consistency loss and 0.2
hard-constraint consistency loss. The released checkpoints correspond to
seeds 31, 47 and 73 and were trained in August 2026. The default
reproduction loads these fixed checkpoints rather than retraining; the
runner exposes retraining separately. Fixed checkpoints do not remove
the decoding runtime boundary: categorical output still depends on the
locked Torch and scikit-learn stack. Same-stack reproduction is the
declared reference; no independent cross-stack GeoFM-LDN raster
comparison is claimed, and changes between pre-lock and
locked-environment outputs are not attributed to CPU architecture or to
one library. During allocation, GeoFM-LDN directly calls the Geospatial
Kernel \texttt{allocate\_action} routine, so both methods share the same
constrained projection and differ primarily in their proposal model. The
comparison is therefore a proposal/pipeline comparison, not a comparison
of two independent projection architectures. The name describes this
benchmark implementation and does not attribute an external publication
or performance claim.

The common evaluator computes actual class counts, normalized demand
total variation, hard-mask violations, strict multi-class change FoM,
binary change FoM, change F1, overall accuracy and macro-F1 for
historical runs. The strict FoM counts a hit only when the predicted
destination class equals the observed destination class; wrong
destination changes are retained in the denominator. It also generates
persistence and minimum-change zero controls, spatial-block bootstrap
intervals and paired model-difference intervals for strict FoM, overall
accuracy and macro-F1. For planning runs it additionally computes
ecological conversion, new-built neighbourhood fraction, new-built
component counts, the union of 2024 built and newly built components,
all-final-built component counts, a 500-m leapfrog rate, distances to
roads and prior built cells, infrastructure proxy distance and built
retirement. Pareto membership is calculated within each scenario over
three release objectives: major-road distance, prior-built distance and
union-built component density. Ecological conversion is retained as a
descriptive public-data pressure proxy because it is structurally zero
for the two exact-count allocators and therefore has no discriminating
power in their comparison. Vegetation gain is a scenario input and
remains descriptive, not an objective.

\subsubsection{Scenario actions and
uncertainty}\label{scenario-actions-and-uncertainty}

The moderate-growth, green-priority-growth and high-outward-growth
actions are stored in the tracked
\texttt{planning\_scenarios\_public\_2025\_2031.json} manifest; the
legacy \texttt{compact} path identifier is retained only for artifact
compatibility. It is not an endorsement of compact-city outcomes, which
can have distributional trade-offs (Burton, 2000). Each action supplies
feasible class totals for 2025--2031. The 2031 totals are an explicit
extension of the previous annual difference and are not derived from an
official population, housing or infrastructure forecast. Future
exogenous rasters are held at their 2024 values. The green-priority
action is not a water-budget model. Each candidate is run at seeds 31,
47 and 73, and planning tables report arithmetic means and population
standard deviations. No p-values are reported because the three seeds
describe computational variability, not independent field samples. The
primary Pareto comparison is within each scenario, contrasting the three
models under the same action; a global nine-candidate frontier is
retained only as lineage. Sensitivity variants remove or add objective
dimensions on the same released rasters. The objective set evolved
through four review-driven revisions and was not preregistered: v1 mixed
feasibility, action and outcome terms; v2 was not released; v3 used
structurally problematic outcomes; and v4 still included
scenario-supplied vegetation gain and all-built fragmentation. The
current release is v6: road access, prior-built distance and union-built
component density. Ecological conversion is reported as a diagnostic
rather than a Pareto objective because it is structurally zero for the
exact-count Kernel and GeoFM-LDN allocations. We treat v6 as a release
objective set rather than a prespecified planning preference, and treat
its sensitivity results as robustness diagnostics.

\subsubsection{Change polygons and
reproducibility}\label{change-polygons-and-reproducibility}

For each model--scenario pair, annual and cumulative differences from
the 2024 raster are extracted as changed 100-m cells and dissolved into
GeoPackage layers. The public delivery manifest covers 45 ensemble
rasters for 2027--2031 and nine vector packages. The audit script
reports \texttt{INCOMPLETE\_INPUTS} or \texttt{FAIL} when source
rasters, reports or vectors are absent; with the released public bundle,
the input audit passed all alignment gates (with a Dynamic World
label-noise warning) and the output audit passed 312 historical,
rolling-origin and planning predictions plus two frozen WorldCover
evidence rasters, with zero failures. The vector layers preserve
raster-cell geometry and should not be interpreted as cadastral parcels.

\subsection{Results}\label{results}

\subsubsection{A common benchmark separates allocation skill from demand
assumptions}\label{a-common-benchmark-separates-allocation-skill-from-demand-assumptions}

The benchmark covers the Abu Dhabi city polygon represented by
OpenStreetMap relation R4479763. The canonical grid is a 475 × 360
lattice in EPSG:32640 at 100-m resolution. After the common coverage and
alignment checks, 79,726 cells were used for scoring. Each cell
represents 1 ha. Annual Dynamic World observations from 2017 to 2024
were harmonized to six land-cover classes: water, woody vegetation, low
vegetation, wetland, built and bare. AlphaEarth annual embeddings, VIIRS
night-time lights, Copernicus DEM derivatives and OpenStreetMap road
distances supplied transition drivers. ESA WorldCover and public
OpenStreetMap geometries supplied water, wetland and infrastructure
exclusion proxies.

The comparison shared the canonical grid, origin states, target actions,
hard masks and evaluator, but the public-data run was not
information-balanced. The original FLUS-style control used seven
continuous drivers in its external ANN suitability model, whereas
Geospatial Kernel used 25 features including current-class indicators
and neighbourhood proportions; GeoFM-LDN used AlphaEarth latent
embeddings and the Kernel allocator. We therefore report the original
result as an unmatched-input pipeline comparison, not evidence that one
learner is intrinsically superior. The additional 25-feature FLUS run is
a feature-space diagnostic only, because it still learns same-year
suitability rather than next-state transitions and does not use the
Kernel projection implementation. Across six platform--seed runs, five
produced zero-change outputs through current-class identity leakage. The
archived FLUS console is a Mach-O arm64 executable rebuilt from the
public GeoSOS source base with the modified
\texttt{train}/\texttt{train-update} entry points and
\texttt{FLUS\_RANDOM\_SEED} seeding patch; its exact source is archived
at an anonymized source-repository revision. It is not asserted to be
bitwise equivalent to an unmodified upstream build. All headline and
diagnostic runs used seeds 31, 47 and 73.

The execution trace makes the distinction between proposal and admission
explicit. A Kernel step records the source state, action,
probability-cube proposal, projected state, model identity, input
evidence and next-state reference. The runtime checks that the action
advances the source time and that the projected raster becomes the next
state. This is the operational meaning of the Kernel in this study: not
a claim that all GWM domains share the same learner, but a claim that
they can share an auditable execution contract. The shared inputs and
execution boundary are summarized in Fig. 1.

\begin{figure}
\centering
\includegraphics[width=1\linewidth,height=\textheight,keepaspectratio,alt={Benchmark and Geospatial Kernel execution contract. Public annual land-cover, embedding, night-light, terrain and road layers are aligned to a common 100-m Abu Dhabi grid. The Kernel exposes a learned proposal, a constraint projection and a state writeback. The schematic describes a public-data benchmark and does not imply statutory planning boundaries.}]{figures/fig01_benchmark_contract.png}
\caption{Benchmark and Geospatial Kernel execution contract. Public
annual land-cover, embedding, night-light, terrain and road layers are
aligned to a common 100-m Abu Dhabi grid. The Kernel exposes a learned
proposal, a constraint projection and a state writeback. The schematic
describes a public-data benchmark and does not imply statutory planning
boundaries.}
\end{figure}

\FloatBarrier

\subsubsection{Historical and planning
results}\label{historical-and-planning-results}

The historical track uses observed target class totals to isolate
spatial allocation from demand estimation. The primary metric is strict
multi-class change FoM: a hit requires the predicted destination class
to equal the observed destination class, while wrong destination changes
remain in the denominator (Pontius et al., 2008). Overall accuracy,
macro-F1, change F1 and binary FoM are secondary diagnostics.
Uncertainty is estimated with an 8 × 8-pixel spatial-block bootstrap,
and model contrasts are computed from the same resampled blocks for each
frozen seed. Persistence and random minimum-change allocation are
explicit zero-model controls.

The planning track starts from the observed 2024 state and applies three
synthetic class-count actions through 2031. The release objective set is
version 6. It uses mean distance of new built cells to major roads, mean
distance to pre-existing built cells and connected components in the
union of 2024 built and newly built cells per 1,000 valid cells. These
are public-data proxies. Ecological conversion is retained as a
descriptive diagnostic, not an objective, because it is structurally
zero for the two exact-count allocators. Vegetation gain is excluded
because it is supplied directly by the scenario action, and all-final-
built component density is reported only as a diagnostic because it can
be changed by built-cell retirement. No weighted composite score is
used. Because the objective set changed across the review rounds, we
disclose the lineage: v1 included seven access, ecological,
neighbourhood and gain terms; v2 proposed roads, prior-built distance,
fragmentation and leapfrog rate but was not run; v3 used ecological
conversion, roads, leapfrog rate and built retirement; and v4 combined
roads, prior-built distance, all-built fragmentation and vegetation
gain. The current v6 release replaces v5 after the review identified the
structural-zero ecological objective and the semantic inversion of
new-only fragmentation. Sensitivity sets are recomputed on the same
rasters and reported as robustness analyses, not as additional
preregistered claims. These are measurable public-data proxies rather
than claims about welfare, equity or statutory zoning.

The released public-data bundle supports a complete re-scoring of the
current historical rasters and a re-compilation of the 2025--2031
planning rasters. The strict multi-class transition FoM (mean across
three seeds) was 0.1256, 0.1961 and 0.1619 for the FLUS-style console,
Geospatial Kernel and GeoFM-LDN in 2023, and 0.1782, 0.2529 and 0.2708
in 2024, respectively. Under this unmatched public-data pipeline,
Geospatial Kernel had the largest 2023 value, whereas GeoFM-LDN had the
largest 2024 value after recursive two-step rollout; no single model
dominated both horizons. These rankings are conditional on the different
inputs, learners and shared evaluation projection.

\textbf{Table 1. Historical strict destination-change FoM.} Values are
means across three frozen computational seeds. The 25-feature FLUS run
is excluded because it is not a comparable estimator: five of six
platform--seed runs collapsed to zero change. Persistence (0.0000 in
both years) and random minimum-change allocation (0.0253 in 2023; 0.0605
in 2024) are explicit zero-model controls and are plotted in Fig. 2.

{\def\LTcaptype{none} % do not increment counter
\begin{longtable}[]{@{}lrrr@{}}
\toprule\noalign{}
Target year & FLUS-style ANN-CA & Geospatial Kernel & GeoFM-LDN \\
\midrule\noalign{}
\endhead
\bottomrule\noalign{}
\endlastfoot
2023 (one-step) & 0.1256 & \textbf{0.1961} & 0.1619 \\
2024 (two-step open-loop) & 0.1782 & 0.2529 & \textbf{0.2708} \\
\end{longtable}
}

Spatial-block bootstrap model contrasts support these differences. For
2023, Kernel minus the FLUS-style control was 0.0703 {[}0.0558,
0.0866{]} and GeoFM-LDN minus Kernel was -0.0341 {[}-0.0479, -0.0209{]}.
For 2024, GeoFM-LDN minus Kernel was 0.0177 {[}0.0033, 0.0314{]}. These
are conditional pipeline contrasts on public labels, not
independent-sample significance tests. The persistence control had
higher overall accuracy than every learned candidate (0.9436 in 2023 and
0.8938 in 2024), illustrating why change-specific metrics are necessary.
Fig. 2 shows the four historical metrics alongside the persistence and
random minimum-change controls; main-model bars use three-seed
population standard deviations, whereas the paired spatial-block
intervals above are the uncertainty summaries used for model contrasts.
The 25- and 19-feature FLUS diagnostics are reported in the
machine-readable archive and Supplementary Table S3.

\begin{figure}
\centering
\includegraphics[width=1\linewidth,height=\textheight,keepaspectratio,alt={Historical allocation skill. Strict destination-change FoM is the primary metric; change F1 is shown as a secondary diagnostic. Bars show mean plus or minus population standard deviation across the three frozen seeds. Persistence and random allocation are explicit zero-model controls.}]{figures/fig02_historical_validation.png}
\caption{Historical allocation skill. Strict destination-change FoM is
the primary metric; change F1 is shown as a secondary diagnostic. Bars
show mean plus or minus population standard deviation across the three
frozen seeds. Persistence and random allocation are explicit zero-model
controls.}
\end{figure}

\FloatBarrier

An additional FLUS input-dimension analysis was also completed. The
25-feature run is not a valid point estimator: macOS seeds 47 and 73 and
all three independent external Windows x86\_64 verification runs
produced zero-change outputs, whereas macOS seed 31 was the sole
non-collapsed run. The six-run pattern is therefore platform-sensitive
and is retained as an identity-leakage diagnostic, not as a Kernel
contrast. The 19-feature diagnostic was run for all three macOS seeds.
Its strict FoM ranged from 0.1071 to 0.1617 in 2023 and from 0.0872 to
0.1386 in 2024; predicted changes ranged from 1,702 to 2,417 and from
1,805 to 2,499 cells, respectively, compared with 4,500 and 8,464
observed changes. Demand total variation ranged from 0.0153 to 0.0305 in
2023 and from 0.0615 to 0.0767 in 2024. These results show why matching
feature dimensions cannot make the tasks equivalent: FLUS learns static
suitability \(p(S_t\mid X_t)\), while the Kernel learns next-state
transitions and then projects them against action deficits and hard
constraints.

The confidence sensitivity is a label-quality and selection-effect
diagnostic, not an independent validation set or a model-skill test.
Requiring both the origin and target Dynamic World mean top probability
to exceed 0.5 retains only 25 of 4,500 observed changes (0.56\%) in 2023
and 56 of 8,464 (0.66\%) in 2024. Strict FoM is 0.000 for all three
models in 2023; in 2024 the means are 0.0040 for the FLUS-style console,
0.0055 for Geospatial Kernel and 0.0188 for GeoFM-LDN. Screening only
the confidence of the year immediately preceding the target avoids
filtering on target-year confidence, but still retains only 129 (2.87\%)
and 93 (1.10\%) observed changes in 2023 and 2024, respectively. Its
2023 FoM values remain 0.000; its 2024 means are 0.0040, 0.0041 and
0.0120 for the FLUS-style console, Kernel and GeoFM-LDN. The two filters
therefore describe label selection rather than establishing or rescuing
model skill. Dynamic World also reports built pixels increasing from
9,359 in 2022 to 15,598 in 2024 (67\%), while the median one-year
reversion fraction is 0.3611 and the mean fraction of cells below
confidence 0.5 is 0.5497. Test-year change counts are approximately
twice the training-transition counts. Because the same grid cells recur
across years and the Kernel includes coordinate features, a
training--test memory effect cannot be excluded; the reported scores
should therefore be interpreted as conditional agreement with the public
product.

\subsubsection{Rolling-origin and external-product tests expose temporal
and label
dependence}\label{rolling-origin-and-external-product-tests-expose-temporal-and-label-dependence}

The leakage-controlled expanding-window test produced mean Kernel strict
FoM values of 0.2595, 0.0724, 0.2116 and 0.1816 for the 2021, 2022, 2023
and 2024 one-step targets, respectively. The corresponding random
minimum-change values were 0.0524, 0.0048, 0.0200 and 0.0323;
persistence produced no predicted change and therefore zero strict FoM.
Paired 8 × 8-cell spatial-block contrasts between Kernel and each zero
model excluded zero in every target year and each of the three seeds.
The shared low 2022 values (Kernel 0.0724; random 0.0048) indicate a
difficult 2021→2022 transition for all allocators, consistent with the
elevated label-reversion rate in Fig. S2 and the contemporaneous
fluctuation in Dynamic World built counts, rather than a Kernel-specific
failure. Persistence also retained higher overall accuracy than Kernel
in 2022--2024, whereas Kernel had higher overall accuracy only in 2021.
The change-specific and all-cell metrics therefore answer different
questions.

The WorldCover comparison has two distinct interpretations. Across built
fractions from 0.05 to 0.30, WorldCover contained 1,445--1,802
built-gain cells and 1,408--1,506 built-loss cells. Dynamic
World-to-WorldCover built-gain F1 was only 0.0111--0.0192. Dynamic World
built pixels decreased from 9,642 in 2020 to 8,553 in 2021, so under the
2020→2021 Dynamic World oracle-demand action the Kernel, random and
persistence allocators necessarily produced no built-gain cells. Their
F1 of 0 is therefore a structural-zero, zero-power result and cannot
evaluate gain-location skill. The informative count-controlled variant
used the WorldCover gain count as the action and ranked 2020-origin
Kernel proposals against a random allocation from the same candidate
population. Kernel F1 was 0.0157--0.0397 versus 0.0251--0.0277 for
random allocation. The mean ± sample-SD pairs (Kernel versus random)
were 0.0157 ± 0.0008 versus 0.0268 ± 0.0058 at 0.05, 0.0228 ± 0.0007
versus 0.0277 ± 0.0036 at 0.10, 0.0268 ± 0.0011 versus 0.0251 ± 0.0011
at 0.20, and 0.0397 ± 0.0007 versus 0.0251 ± 0.0075 at 0.30. Thus Kernel
was below random at 0.05, the 0.10--0.20 differences were not
distinguishable from three-seed variation, and only the 0.30 result was
clearly above random. The low-threshold deficit is consistent with
WorldCover cells that barely exceed 5\% built fraction, including sparse
or road-edge structure that Dynamic World does not consistently label as
built; the Kernel proposal learned from Dynamic World therefore ranks
these product- specific gains later than random candidates. For 2021
built stock, Kernel F1 was 0.3518--0.3971, while persistence was
0.3886--0.4387. The corresponding Dynamic World 2020 and 2021
observed-stock baselines were 0.3908--0.4334 and 0.3501--0.3924,
respectively, showing that the apparent persistence advantage mirrors
greater agreement of the 2020 label with WorldCover rather than a model
advantage. Fig. 3 separates these stock and count-controlled gain
comparisons; none is authoritative validation.

\begin{figure}
\centering
\includegraphics[width=1\linewidth,height=\textheight,keepaspectratio,alt={Rolling-origin robustness and external-product diagnostics. Panel a reports strict destination-change FoM for four leakage-controlled one-step targets; points show means and sample standard deviations across three computational seeds. Panel b reports Kernel-minus-baseline paired spatial-block contrasts; points are mean seed medians and whiskers span the minimum to maximum of seed-specific 95\% intervals from 1,000 resamples of 8 × 8-cell blocks. Panel c compares same-year Dynamic World 2020 and 2021 built stocks, plus 2021 Kernel and persistence predictions, with the corresponding WorldCover products at four 100-m built-fraction thresholds. Panel d compares Dynamic World built gain with the count-controlled Kernel and random allocations, where WorldCover gain count is supplied as the action. The original oracle-demand gain test is structural-zero and is not plotted as model evidence. The WorldCover comparison is an external public-product diagnostic, not authoritative ground truth. Source data are provided in source\_data\_fig03\_rolling\_external\_diagnostics.csv.}]{figures/fig03_rolling_external_diagnostics.png}
\caption{Rolling-origin robustness and external-product diagnostics.
Panel a reports strict destination-change FoM for four
leakage-controlled one-step targets; points show means and sample
standard deviations across three computational seeds. Panel b reports
Kernel-minus-baseline paired spatial-block contrasts; points are mean
seed medians and whiskers span the minimum to maximum of seed-specific
95\% intervals from 1,000 resamples of 8 × 8-cell blocks. Panel c
compares same-year Dynamic World 2020 and 2021 built stocks, plus 2021
Kernel and persistence predictions, with the corresponding WorldCover
products at four 100-m built-fraction thresholds. Panel d compares
Dynamic World built gain with the count-controlled Kernel and random
allocations, where WorldCover gain count is supplied as the action. The
original oracle-demand gain test is structural-zero and is not plotted
as model evidence. The WorldCover comparison is an external
public-product diagnostic, not authoritative ground truth. Source data
are provided in
source\_data\_fig03\_rolling\_external\_diagnostics.csv.}
\end{figure}

\FloatBarrier

The reference reproduction environment is macOS arm64 with Python 3.11
and scikit-learn 1.9.0. The archived Linux arm64 rerun uses Python 3.12
and scikit-learn 1.8.0, so it measures a cross-stack boundary rather
than an architecture-only effect. Across six rasters (2023 and 2024 for
each seed), it differs from the current reference by 498--1,147 cells
(0.62--1.44\% of valid cells), with strict-FoM deltas from -0.0023 to
0.0010. An controlled GitHub Actions Ubuntu x86\_64 rerun under Python
3.11.16 and scikit-learn 1.9.0 matched all six released categorical
reference rasters exactly; its archived audit is linked in Code
availability. This is a scoped locked-stack regression result, not
universal bitwise identity. Independent external Windows FLUS
verification is documented with its provenance limitation in
Supplementary Table S3; its source rasters are not locally archived. The
archived comparison report records the exact paths and environment
fields for the arm64 pair.

The planning compiler compares the three models within each of the three
scenarios. At 2031, the FLUS-style control has union-built component
densities of 2.48--3.27 per 1,000 valid cells, the Kernel 2.92--4.28 and
GeoFM-LDN 3.21--4.35. The Kernel is closer to major roads and prior
built cells, whereas the union metric shows that this edge-filling
pattern does not automatically imply lower overall fragmentation than
the FLUS-style allocations. Under the release objective set, the
within-scenario frontiers contain FLUS-style and Kernel candidates in
all three scenarios; GeoFM-LDN is dominated in each scenario. In the
preceding v5 formulation, GeoFM-LDN could appear non-dominated only
because its ecological-conversion value was the same structural zero as
the exact-count Kernel; that zero is not a model advantage. Ecological
conversion is therefore retained only as a diagnostic in v6. All
sensitivity variants retain the same FLUS-style/Kernel two-model
frontier. These statements are conditional on public proxy layers and
synthetic actions, not model superiority claims. The corresponding
objective profiles and diagnostic trade-offs are shown in Fig. 4.

\begin{figure}
\centering
\includegraphics[width=1\linewidth,height=\textheight,keepaspectratio,alt={Conditional planning objectives. The 2031 outcomes compare the three models under moderate-growth, green-priority-growth and high-outward-growth actions. Panels show the three release objectives: distance to major roads, distance to pre-existing built cells, and connected components in the combined 2024-built and newly built footprint. Black outlines indicate membership of the within-scenario Pareto frontier.}]{figures/fig04_planning_objectives.png}
\caption{Conditional planning objectives. The 2031 outcomes compare the
three models under moderate-growth, green-priority-growth and
high-outward-growth actions. Panels show the three release objectives:
distance to major roads, distance to pre-existing built cells, and
connected components in the combined 2024-built and newly built
footprint. Black outlines indicate membership of the within-scenario
Pareto frontier.}
\end{figure}

\FloatBarrier

\textbf{Table 2. 2031 within-scenario planning objectives.} Values are
means across three seeds. Road and pre-existing-built distances are in
metres; components are the number in the combined 2024-built and newly
built footprint per 1,000 valid cells. Ecological conversion is a
descriptive fraction of new-built cells whose origin was woody
vegetation, low vegetation or wetland. Asterisks indicate the primary
within-scenario Pareto frontier.

{\def\LTcaptype{none} % do not increment counter
\begin{longtable}[]{@{}
  >{\raggedright\arraybackslash}p{(\linewidth - 10\tabcolsep) * \real{0.1600}}
  >{\raggedleft\arraybackslash}p{(\linewidth - 10\tabcolsep) * \real{0.1600}}
  >{\raggedleft\arraybackslash}p{(\linewidth - 10\tabcolsep) * \real{0.1600}}
  >{\raggedleft\arraybackslash}p{(\linewidth - 10\tabcolsep) * \real{0.1600}}
  >{\raggedleft\arraybackslash}p{(\linewidth - 10\tabcolsep) * \real{0.1600}}
  >{\centering\arraybackslash}p{(\linewidth - 10\tabcolsep) * \real{0.2000}}@{}}
\toprule\noalign{}
\begin{minipage}[b]{\linewidth}\raggedright
Scenario / model
\end{minipage} & \begin{minipage}[b]{\linewidth}\raggedleft
Road (m)
\end{minipage} & \begin{minipage}[b]{\linewidth}\raggedleft
Prior built (m)
\end{minipage} & \begin{minipage}[b]{\linewidth}\raggedleft
Components / 1,000
\end{minipage} & \begin{minipage}[b]{\linewidth}\raggedleft
Ecology*
\end{minipage} & \begin{minipage}[b]{\linewidth}\centering
Frontier
\end{minipage} \\
\midrule\noalign{}
\endhead
\bottomrule\noalign{}
\endlastfoot
\textbf{Moderate} & & & & & \\
FLUS-style ANN--CA & 446.1 & 231.6 & 3.035 & 0.0168 & * \\
Geospatial Kernel & 334.7 & 116.7 & 3.943 & 0.0000 & * \\
GeoFM-LDN & 475.9 & 239.7 & 4.056 & 0.0000 & \\
\textbf{Green priority} & & & & & \\
FLUS-style ANN--CA & 443.7 & 212.7 & 3.265 & 0.0177 & * \\
Geospatial Kernel & 323.1 & 111.0 & 4.277 & 0.0000 & * \\
GeoFM-LDN & 452.1 & 222.3 & 4.348 & 0.0000 & \\
\textbf{High outward} & & & & & \\
FLUS-style ANN--CA & 443.5 & 255.3 & 2.484 & 0.0131 & * \\
Geospatial Kernel & 365.5 & 139.8 & 2.923 & 0.0000 & * \\
GeoFM-LDN & 507.5 & 263.0 & 3.207 & 0.0000 & \\
\end{longtable}
}

\subsubsection{Sensitivity shows a stable but not fully explained
neighbourhood
contribution}\label{sensitivity-shows-a-stable-but-not-fully-explained-neighbourhood-contribution}

The base Kernel allocation score adds a 7 × 7 target-class neighbourhood
fraction with weight 0.35. The existing sensitivity and mechanism
artefacts are retained as diagnostic controls; they do not identify a
causal effect because the proposal features and allocator both contain
spatial context. The neighbourhood term is therefore interpreted as
evidence about the execution mechanism, not as an independently
estimated planning preference. Fig. 5A shows focused proposal controls,
and Fig. 5B reports the matched runtime controls.

In the post-hoc neighbourhood-weight sensitivity, mean strict FoM ranged
from 0.19570 to 0.19609 in 2023 and from 0.25273 to 0.25304 in 2024; the
complete strict-evaluator scan is provided in Supplementary Table S2.
These changes are small relative to the model contrasts and are not
interpreted causally. Objective-set sensitivity on the planning rasters
retained the FLUS-style and Kernel candidates in all scenarios.
GeoFM-LDN was dominated in all three scenarios after the structurally
zero ecological-conversion term was removed from the release objectives.

\subsubsection{Mechanism controls separate learned proposal from runtime
projection}\label{mechanism-controls-separate-learned-proposal-from-runtime-projection}

Mechanism controls separate the learned proposal from runtime
projection, hard constraints and recursive writeback. They are not
causal policy experiments. Deleting the action reduces mean strict FoM
by 0.1124 in 2023 and 0.1699 in 2024, whereas deleting the hard
constraint increases it by 0.014 and 0.011, respectively. The former is
expected because the historical task supplies oracle target counts; it
shows that much of the measured skill belongs to the action-conditioned
allocator rather than to a demand forecast. The latter is an important
warning: allowing changes inside the exclusion mask can improve
agreement with noisy labels while violating the planning contract. The
controls therefore support an execution-level interpretation but do not
prove that either component causes an urban outcome.

\begin{figure}
\centering
\includegraphics[width=1\linewidth,height=\textheight,keepaspectratio,alt={Mechanism controls. Panel A compares the full transition proposal with selected proposal controls. Panel B gives the strict destination-change FoM difference between matched runtime controls and the full Kernel. Bars show mean plus or minus population standard deviation across the three frozen seeds. The no-state-writeback control is undefined for the 2023 one-step target and is therefore shown only for 2024. These are execution diagnostics, not causal policy experiments.}]{figures/fig05_mechanism_ablation.png}
\caption{Mechanism controls. Panel A compares the full transition
proposal with selected proposal controls. Panel B gives the strict
destination-change FoM difference between matched runtime controls and
the full Kernel. Bars show mean plus or minus population standard
deviation across the three frozen seeds. The no-state-writeback control
is undefined for the 2023 one-step target and is therefore shown only
for 2024. These are execution diagnostics, not causal policy
experiments.}
\end{figure}

\FloatBarrier

\subsubsection{Spatial footprints are scenario-conditioned, not parcel
predictions}\label{spatial-footprints-are-scenario-conditioned-not-parcel-predictions}

The 2031 majority-vote ensemble L1 demand errors were 8, 44 and 42
pixels for the Kernel's moderate, green-priority and high-outward
scenarios, respectively; the corresponding GeoFM-LDN errors were 1,580,
1,524 and 1,148 pixels and the FLUS-style errors were 2,504, 2,888 and
2,568 pixels. Seed-level projections meet their feasible class totals,
but majority voting can produce a different raster and therefore a
non-zero ensemble error. All 312 historical, rolling-origin and planning
prediction records passed the raster audit, and the two frozen
WorldCover evidence rasters passed the evidence audit. Fig. 6 uses the
high-outward scenario to make the spatial contrast legible; the complete
nine-panel atlas is provided as Fig. S1. These are raster-cell change
footprints, not parcel boundaries or statutory zoning maps.

\begin{figure}
\centering
\includegraphics[width=1\linewidth,height=\textheight,keepaspectratio,alt={High-outward-growth spatial footprints. The observed 2024 state is compared with each model's 2031 majority-vote ensemble. Orange denotes cells newly built between 2024 and 2031; grey denotes cells already built in 2024. The map is a 100-m scenario stress-test product, not a cadastral or statutory land-use map.}]{figures/fig06_planning_maps_2031.png}
\caption{High-outward-growth spatial footprints. The observed 2024 state
is compared with each model's 2031 majority-vote ensemble. Orange
denotes cells newly built between 2024 and 2031; grey denotes cells
already built in 2024. The map is a 100-m scenario stress-test product,
not a cadastral or statutory land-use map.}
\end{figure}

\FloatBarrier

\subsection{Discussion}\label{discussion}

This study positions Geospatial Kernel as a candidate execution layer
for constrained spatial reasoning and, more broadly, proposes an audit
protocol for noisy annual labels. The defensible contribution is that a
transition proposal and a planning admission rule can be represented and
inspected separately. Under the unmatched public-data pipeline,
Geospatial Kernel has the largest strict transition FoM at the one-step
horizon, while GeoFM-LDN has the largest value after recursive two-step
rollout. In the separate rolling analysis, Kernel exceeds both zero
models in all four one-step windows, but its score varies from 0.0724 to
0.2595 after excluding the later road snapshot. The WorldCover
diagnostic does not provide a blanket external validation: the original
Dynamic World-count gain comparison is structurally zero-powered because
the observed built count declines. The count-controlled variant remains
near opportunity-level agreement overall: Kernel is below random at
0.05, indistinguishable from random at 0.10--0.20 given three-seed
variation, and above random only at 0.30. The low-threshold deficit is
consistent with WorldCover's sparse or road-edge cells just above the
5\% built-fraction threshold being inconsistently represented by Dynamic
World. Confidence filtering, rolling-window stability and
external-product agreement are different tests, and none substitutes for
authoritative local validation. We therefore treat the planning frontier
as conditional on public proxy objectives and synthetic actions.

The evidence remains bounded. The six classes are remote-sensing land
cover, not statutory residential, commercial or industrial land use.
Dynamic World labels show substantial annual uncertainty and apparent
turnover; the available 2020--2021 WorldCover diagnostic is neither
fully independent nor authoritative, and no external 2023--2024 change
product or observed 2025--2031 labels was available. The headline
analysis also uses a 31 July 2026 OSM road snapshot, which may contain
roads constructed after 2022. This information is shared by all three
pipelines and therefore does not change their ordering, but it can
affect absolute scores; the strict rolling-origin analysis removes those
road features and should define the no-leakage boundary. As a
descriptive bound, the road-free rolling 2023 one-step FoM is 0.2116, of
the same order as the headline 2023 value of 0.1961; these values are
not directly comparable because the rolling fit uses 23 features and
training through 2022, whereas the headline fit uses 25 features and
training through 2021. Public wetland and protected-area layers are
proxies, future drivers are held at 2024, and the green-priority action
has no water-budget constraint. The FLUS-style control is an
Apple-Silicon binary rebuilt from the public GeoSOS source base with
\texttt{train}/\texttt{train-update} and deterministic-seeding patches;
its headline comparison uses seven drivers. The 25-feature run is a
platform-sensitive identity-leakage diagnostic: five of six
platform--seed runs collapsed to zero change, so the non-collapsed macOS
seed-31 result is not used as a comparative estimate. These public-data
stress tests therefore do not establish causal planning effects or a
validated Abu Dhabi forecast. An operational deployment would replace
public labels and masks with locally approved land-use, infrastructure,
ecological and development-demand records while retaining the same
execution contract.

\subsection{Data availability}\label{data-availability}

The public-data benchmark manifests, protocols, reports and generated
delivery artefacts are available from an anonymized repository and
archival package. Repository and archive identifiers are withheld from
this manuscript during double-anonymized review and will be supplied in
the accepted version. The released package contains the aligned
2017--2024 public rasters, historical and rolling-origin prediction
rasters, frozen WorldCover evidence rasters, 2025--2031 planning
rasters, dissolved raster-cell transition footprints, model assets and
integrity manifests. Authoritative Abu Dhabi land-use and independent
reference data are not part of this study.

\subsection{Code availability}\label{code-availability}

The analysis scripts, benchmark protocol, model runners,
figure-generation code, model assets and integrity checks are maintained
in the accompanying anonymized repository and archival package.
Identifiers and contact details are withheld during double-anonymized
review and will be supplied in the accepted version.

\subsection{Declarations}\label{declarations}

Funding, competing interests, author contributions, acknowledgements and
ethics information are supplied in the separate title page and
submission metadata.

\subsection{References}\label{references}

Abu Dhabi Urban Planning Council. (2007). Plan Abu Dhabi 2030: urban
structure framework plan. Abu Dhabi, United Arab Emirates.

Brown, C. F., Brumby, S. P., Guzder-Williams, B., Birch, T., Hyde, S.
B., Mazzariello, J., Czerwinski, W., Pasquarella, V. J., Haertel, R.,
Ilyushchenko, S., Schwehr, K., Weisse, M., Stolle, F., Hanson, C.,
Guinan, O., Moore, R., \& Tait, A. M. (2022). Dynamic World, near
real-time global 10 m land use land cover mapping. \emph{Scientific
Data}, 9, 251. https://doi.org/10.1038/s41597-022-01307-4

Brown, C. F., Kazmierski, M. R., Pasquarella, V. J., Rucklidge, W. J.,
Samsikova, M., Zhang, C., Shelhamer, E., Lahera, E., Wiles, O.,
Ilyushchenko, S., Gorelick, N., Zhang, L. L., Alj, S., Schechter, E.,
Askay, S., Guinan, O., Moore, R., Boukouvalas, A., \& Kohli, P. (2025).
AlphaEarth Foundations: An embedding field model for accurate and
efficient global mapping from sparse label data. \emph{arXiv}.
https://doi.org/10.48550/arXiv.2507.22291

Burton, E. (2000). The compact city: just or just compact? A preliminary
analysis. \emph{Urban Studies}, 37, 1969--2006.
https://doi.org/10.1080/00420980050162184

Elvidge, C. D., Zhizhin, M., Ghosh, T., Hsu, F.-C., \& Taneja, J.
(2021). Annual time series of global VIIRS nighttime lights derived from
monthly averages: 2012 to 2019. \emph{Remote Sensing}, 13, 922.
https://doi.org/10.3390/rs13050922

European Space Agency. (2024). Copernicus DEM GLO-30: global 30 m
digital elevation model {[}Data set{]}. Google Earth Engine Data
Catalog.
https://developers.google.com/earth-engine/datasets/catalog/COPERNICUS\_DEM\_GLO30\_2024\_1

GeoSOS team, Sun Yat-sen University. (n.d.). FLUS source code for
non-commercial academic research.
http://www.geosimulation.cn/FLUS-source-code.html (accessed 6 September
2026).

Gorelick, N., Hancher, M., Dixon, M., Ilyushchenko, S., Thau, D., \&
Moore, R. (2017). Google Earth Engine: planetary-scale geospatial
analysis for everyone. \emph{Remote Sensing of Environment}, 202,
18--27. https://doi.org/10.1016/j.rse.2017.06.031

Lau, K. H., \& Kam, B. H. (2005). A cellular automata model for urban
land-use simulation. \emph{Environment and Planning B: Planning and
Design}, 32, 247--263. https://doi.org/10.1068/b31110

Liang, X., Guan, Q., Clarke, K. C., Liu, S., Wang, B., \& Yao, Y.
(2021). Understanding the drivers of sustainable land expansion using a
patch-generating land use simulation (PLUS) model: A case study in
Wuhan, China. \emph{Computers, Environment and Urban Systems}, 85,
101569. https://doi.org/10.1016/j.compenvurbsys.2020.101569

Liu, X., Liang, X., Li, X., Xu, X., Ou, J., Chen, Y., Li, S., Wang, S.,
\& Pei, F. (2017). A future land use simulation model (FLUS) for
simulating multiple land use scenarios by coupling human and natural
effects. \emph{Landscape and Urban Planning}, 168, 94--116.
https://doi.org/10.1016/j.landurbplan.2017.09.019

OpenStreetMap contributors. (2026). OpenStreetMap data, Abu Dhabi
relation R4479763 and Geofabrik GCC extract, accessed 31 July 2026.
https://www.openstreetmap.org

Pedregosa, F., Varoquaux, G., Gramfort, A., Michel, V., Thirion, B.,
Grisel, O., Blondel, M., Prettenhofer, P., Weiss, R., Dubourg, V.,
Vanderplas, J., Passos, A., Cournapeau, D., Brucher, M., Perrot, M., \&
Duchesnay, E. (2011). Scikit-learn: Machine learning in Python.
\emph{Journal of Machine Learning Research}, 12, 2825--2830.
https://jmlr.org/papers/v12/pedregosa11a.html

Pontius, R. G., Jr., Boersma, W., Castella, J.-C., Clarke, K., de Nijs,
T., Dietzel, C., Duan, Z., Fotsing, E., Goldstein, N., Kok, K., Koomen,
E., Lippitt, C. D., McConnell, W., Mohd Sood, A., Pijanowski, B.,
Pithadia, S., Sweeney, S., Trung, T. N., Veldkamp, A. T., \& Verburg, P.
H. (2008). Comparing the input, output, and validation maps for several
models of land change. \emph{The Annals of Regional Science}, 42,
11--37. https://doi.org/10.1007/s00168-007-0138-2

Pontius, R. G., Jr., \& Millones, M. (2011). Death to Kappa: birth of
quantity disagreement and allocation disagreement for accuracy
assessment. \emph{International Journal of Remote Sensing}, 32,
4407--4429. https://doi.org/10.1080/01431161.2011.552923

Silva, E. A., \& Clarke, K. C. (2002). Calibration of the SLEUTH urban
growth model for Lisbon and Porto, Portugal. \emph{Computers,
Environment and Urban Systems}, 26, 525--552.
https://doi.org/10.1016/S0198-9715(01)00014-X

Soares-Filho, B. S., Cerqueira, G. C., \& Pennachin, C. L. (2002).
DINAMICA---a stochastic cellular automata model designed to simulate the
landscape dynamics in an Amazonian colonization frontier.
\emph{Ecological Modelling}, 154, 217--235.
https://doi.org/10.1016/S0304-3800(02)00059-5

van Vliet, J., Bregt, A. K., \& Hagen-Zanker, A. (2011). Revisiting
Kappa to account for change in the accuracy assessment of land-use
change models. \emph{Ecological Modelling}, 222, 1367--1375.
https://doi.org/10.1016/j.ecolmodel.2011.01.017

van Vliet, J., Bregt, A. K., Brown, D. G., van Delden, H., Heckbert, S.,
\& Verburg, P. H. (2016). A review of current calibration and validation
practices in land-change modeling. \emph{Environmental Modelling \&
Software}, 82, 174--182. https://doi.org/10.1016/j.envsoft.2016.04.017

Verburg, P. H., Soepboer, W., Veldkamp, A., Limpiada, R., Espaldon, V.,
\& Mastura, S. S. A. (2002). Modeling the spatial dynamics of regional
land use: the CLUE-S model. \emph{Environmental Management}, 30,
391--405. https://doi.org/10.1007/s00267-002-2630-x

Zanaga, D., Van De Kerchove, R., De Keersmaecker, W., Souverijns, N.,
Brockmann, C., Quast, R., Wevers, J., Grosu, A., Paccini, A., Vergnaud,
S., Cartus, O., Santoro, M., Fritz, S., Georgieva, I., Lesiv, M.,
Carter, S., Herold, M., Li, L., Tsendbazar, N.-E., Ramoino, F., \&
Arino, O. (2021). ESA WorldCover 10 m 2020 v100 {[}Data set{]}. Zenodo.
https://doi.org/10.5281/zenodo.5571936

Zanaga, D., Van De Kerchove, R., Daems, D., De Keersmaecker, W.,
Brockmann, C., Kirches, G., Wevers, J., Cartus, O., Santoro, M., Fritz,
S., Lesiv, M., Herold, M., Tsendbazar, N.-E., Xu, P., Ramoino, F., \&
Arino, O. (2022). ESA WorldCover 10 m 2021 v200 {[}Data set{]}. Zenodo.
https://doi.org/10.5281/zenodo.7254221

\subsection{Supplementary information
outline}\label{supplementary-information-outline}

\begin{itemize}
\tightlist
\item
  \textbf{Supplementary Note 1:} Complete data crosswalk, quality
  metrics and public/proxy status.
\item
  \textbf{Supplementary Note 2:} Kernel runtime schemas, audit fields
  and adapter boundary.
\item
  \textbf{Supplementary Table S1:} Full 2025--2031 model--scenario
  objective matrix with per-seed values.
\item
  \textbf{Supplementary Table S2:} Neighbourhood-weight sensitivity at
  0, 0.175, 0.35 and 0.7 (see
  \texttt{supplementary\_table\_S2\_neighbourhood\_weight\_sensitivity.md}).
\item
  \textbf{Supplementary Table S3:} FLUS 13-, 19- and 25-feature
  diagnostics, platform boundary and demand-underfill evidence (see
  \texttt{supplementary\_table\_S3\_flus\_feature\_diagnostics.md}).
\item
  \textbf{Supplementary Table S4:} Raster and vector delivery audit,
  hashes and layer counts.
\item
  \textbf{Supplementary Table S5:} Rolling-origin and WorldCover
  external-product diagnostics (see
  \texttt{supplementary\_table\_S5\_rolling\_external\_diagnostics.md}).
\item
  \textbf{Fig. S1:} Full 3-scenarios × 3-model 2031 new-built footprint
  atlas (\texttt{figures/figS01\_planning\_atlas\_2031.*}).
\item
  \textbf{Fig. S2:} Dynamic World confidence and high-confidence change
  sensitivity (\texttt{figures/figS02\_input\_label\_quality.*}).
\item
  \textbf{Fig. S3:} Historical maps and change-error maps for 2024
  (\texttt{figures/figS03\_historical\_2024\_maps\_and\_errors.*}).
\item
  \textbf{Fig. S4:} Driver layers and experiment design in English
  (\texttt{figures/figS04\_driver\_layers\_and\_experiment\_design.*}).
\end{itemize}

\begin{center}\rule{0.5\linewidth}{0.5pt}\end{center}

\end{document}
